% Wayfinder manuscript — master document (LightsOut LaTeX approach, journal-adapted)
% Compile: latexmk -pdf main.tex   (or .\compile.ps1)
% Sections live in sections/*.tex (generated from the markdown by build_tex.py); edit prose there
% or in the markdown source and re-run build_tex.py.

\documentclass[11pt]{article}

% --- encoding / fonts (helvet sans, the LightsOut look) ---
\usepackage[utf8]{inputenc}
\DeclareUnicodeCharacter{03C9}{\ensuremath{\omega}}  % omega in a bib title
\DeclareUnicodeCharacter{2014}{---}                  % em-dash in a bib field
\usepackage[T1]{fontenc}
\usepackage{helvet}
\renewcommand{\familydefault}{\sfdefault}
\usepackage[margin=1in]{geometry}
\usepackage{setspace}\onehalfspacing
\setlength{\parindent}{0pt}\setlength{\parskip}{4pt}
\tolerance=400\emergencystretch=3pt

% --- colours + section headers (LightsOut aesthetic; starred = unnumbered so the manual
%     "4.1b"-style numbers in titles match the prose's literal cross-references) ---
\usepackage[dvipsnames]{xcolor}
\definecolor{sectioncolor}{RGB}{0,0,0}      % black (standard manuscript headings)
\definecolor{subsectioncolor}{RGB}{0,0,0}
\definecolor{markercolor}{RGB}{0,0,0}
\usepackage{titlesec}
\titleformat{\section}{\normalfont\bfseries\large}{}{0em}{}
  [\vspace{-0.4em}\rule{\textwidth}{0.4pt}]
\titleformat{\subsection}{\normalfont\bfseries}{}{0em}{}
\titleformat{\subsubsection}[runin]{\normalfont\bfseries\itshape}{}{0em}{}[:]
\titlespacing*{\section}{0pt}{10pt}{4pt}
\titlespacing*{\subsection}{0pt}{6pt}{2pt}
\titlespacing*{\subsubsection}{0pt}{4pt}{0.5em}

% --- math / tables / figures / links ---
\usepackage{amsmath}
\usepackage{calc}  % defines \real{} for pandoc's longtable column-width arithmetic
\usepackage{graphicx,booktabs,longtable,array,caption}
\captionsetup{font=footnotesize,labelfont=bf,skip=3pt}
\usepackage[hidelinks]{hyperref}

% --- references (numbered superscript, LightsOut style; swap to author-year for Frontiers final) ---
\usepackage[numbers,super,sort&compress]{natbib}
\bibliographystyle{unsrtnat}

% --- pandoc helper macros (fragments use these) ---
\providecommand{\tightlist}{\setlength{\itemsep}{0pt}\setlength{\parskip}{0pt}}
\providecommand{\passthrough}[1]{#1}
\providecommand{\pandocbounded}[1]{#1}
% pandoc longtables without captions + hyperref reference a counter literally named 'none'
\providecommand{\theHnone}{none.\arabic{none}}
\newcounter{none}

% --- title ---
\title{\vspace{-2em}\bfseries\color{sectioncolor}Receipt-backed prioritization for literature-based\\
  discovery using Perturb-seq evidence}
\author{Dayanjan Shanaka Wijesinghe}
\date{\today}

\begin{document}
\maketitle

\begin{abstract}
\noindent
Literature-based discovery (LBD) connects facts that are individually established but never joined, and
four decades of work have made it a prolific generator of hypotheses. Its persistent weakness is not
generation but triage: an LBD search returns far more candidate links than any laboratory can pursue, and
ranking by literature-rarity tracks obscurity as readily as importance.

\vspace{2pt}\noindent
We present \textbf{Wayfinder}, an approach that addresses this gap by pairing LBD generation with a
deterministic referee that evaluates each proposed \emph{gene $\rightarrow$ T-cell program $\rightarrow$
disease} hypothesis against a held experimental substrate --- a genome-scale CRISPRi Perturb-seq resource in
primary human CD4$^{+}$ T cells --- and returns, for each, a prioritization verdict with a receipt at every
hop. We keep two receipt classes distinct: the knockdown, effect, and program hops are backed by
experimental perturbation measurements, whereas the disease hop is a genetic-association nomination rather
than a claim of experimental causality. Two diagnostics define the referee's behaviour. It is willing to
answer a confident, receipt-backed \emph{no} --- a \emph{falsification} diagnostic --- and it distinguishes a
\emph{refuted} hypothesis from an \emph{untested} one, an explicit \emph{abstention}: a failed knockdown is
reported as untested, an artifact caught rather than a false negative recorded.

\vspace{2pt}\noindent
The entire loop --- generation, evidence integration, and provenance --- ran inside an agentic scientific
workbench in which an independent critic model enforced calibrated \emph{language} on the platform's own
output, and its analyses reproduced when the workbench was driven programmatically --- the pipeline's funnel
and ranking digit-for-digit, and the sensitivity panel byte-for-byte (replicable in principle, subject to
the platform's user interface). Applied to the CD4$^{+}$ resource, the approach culled 22{,}039
machine-generated hypotheses to a small set of receipt-backed survivors; a verdict ledger demonstrates that
the referee discriminates, and negative controls locate that discrimination honestly --- including the
finding that the disease hop's stringency is inherited from the substrate rather than supplied by the
referee. The worked example --- NAB2 $\rightarrow$ Th1/Th2 polarization $\rightarrow$ atopic eczema --- is a
literature-novel regulatory nomination whose sharpest artifactual confounder, a CRISPRi \emph{cis}-effect on
the adjacent \emph{STAT6}, we falsify against the study authors' own genome-wide data, and which an
independent cross-model replication reproduced. The contribution is the method and what it found:
\textbf{receipt-backed prioritization of machine-generated hypotheses, with explicit abstention and
falsification diagnostics} --- not a demonstration of predictive correctness.

\vspace{4pt}\noindent\textbf{Keywords:} literature-based discovery; hypothesis prioritization; Perturb-seq;
CRISPRi screens; T-cell polarization; agentic language models; falsification; abstention.
\end{abstract}

\section*{1. Introduction}A researcher today works between two floods. On one side is the literature: more than a million new
articles appear each year, and any single mechanistic question is now spread across more papers than one
person can read. On the other side is the data: a single genome-scale screen returns effect estimates for
thousands of genes across thousands of genetic perturbations, and a lab that runs one is often left
holding a matrix of measurements with no obvious question to ask of it. The two floods share a shape ---
far more has been \emph{measured} and \emph{written down} than has been \emph{connected} --- and the gap between them is
where testable connections wait, hidden in plain sight.

Literature-based discovery (LBD) was designed to bridge exactly that gap. Since Swanson's demonstration
that fish oil and Raynaud's disease were linked by facts already in print but never joined \cite{swanson1986},
LBD has framed a hypothesis as a triangle: an entity \textbf{A} relates to an intermediate program \textbf{B}, \textbf{B}
relates to an outcome \textbf{C}, yet \textbf{A} and \textbf{C} have never been connected directly. If both legs are
supported and the closing edge is missing from the literature, that missing edge is a candidate
connection --- novel in the literature, not yet a claim about biology.
The paradigm is now four decades old and has been applied across biomedicine, our own prior work included:
we used an ABC co-occurrence pipeline to connect a plasma-metabolite signal to a druggable target in
cardiac arrest \cite{henry2021}, a link for which we then found evidence in vivo \cite{cheng2021}.

But LBD has a chronic, well-documented weakness, and it is worth stating in the words that named it. In
the review of a proposal submitted for funding consideration, three reviewers
converged on a single objection:

\begin{quote}
\emph{``This application has the same major problem that has plagued all LBD work: it generates an enormous
number of hypotheses, almost none of which ever get followed up.''}
\end{quote}

They were right. LBD is a generator, and a generator with no discriminating filter overproduces: it turns
one flood into another. Classic LBD ranks candidates by properties of the \emph{literature} --- how rare a
co-mention is, how strong an intermediate link looks --- but rarity in the literature tracks \emph{obscurity} as
readily as it tracks \emph{importance}, and a ranked list, however long, is not an answer. What has been
missing is not more generation but adjudication: a way to ask, of each machine-generated hypothesis,
\emph{is this actually supported by data?} --- and to be willing to answer \textbf{no}.

Here we describe \textbf{Wayfinder}, an approach to closing that loop. The approach pairs LBD generation with
a deterministic referee that adjudicates each proposed \emph{gene $\rightarrow$ T-cell program $\rightarrow$ disease} triple against a
held experimental substrate --- a genome-scale CRISPRi Perturb-seq resource in primary human CD4+ T cells
\cite{zhu2025} --- and returns a verdict with a receipt at every hop. We instantiate the
approach in a concrete, reproducible implementation, but the contribution is the \emph{method} --- receipt-backed
adjudication of machine-generated hypotheses --- and what it found, not a software product.
The receipts are of two kinds, and we keep them distinct: the knockdown, effect, and program hops are
backed by \emph{experimental} measurements from the perturbation data; the disease hop is an \emph{association}
receipt (genetic enrichment), a nomination rather than a claim of experimental causality.

Two features define the approach. First, the referee is willing to say a confident, receipt-backed \textbf{no},
and it distinguishes a \emph{refuted} hypothesis from an \emph{untested} one: when a gene's knockdown fails quality
control, the result is reported as \emph{untested}, never as a negative --- an artifact caught rather than a false
negative recorded. Falsification, not confirmation, is the point. Second, the entire loop --- generation,
adjudication, and the assembly of its own provenance --- was run inside an agentic scientific workbench
(Claude Science), in which one model authored the analysis and an independent reviewer model, at separate
checkpoints, verified every number and enforced calibrated language on the manuscript-facing output,
flagging and removing overstated words from the platform's own text. The platform audited its own output ---
through an independent reviewer \emph{role} (a distinct model at distinct checkpoints), not a cross-vendor
check; cross-family independence is provided separately (Section 4.6).

We demonstrate the Wayfinder approach on the CD4+ T-cell resource, where it posed 22,039
gene$\rightarrow$program$\rightarrow$disease hypotheses that a literature-novelty gate and a deterministic referee together culled
to a small set of receipt-backed survivors. A ledger of adjudicated verdicts \emph{demonstrates} that the
referee discriminates --- supporting, refuting, and declining to test --- and a negative-control panel ---
failed-knockdown genes, which must return \emph{untested}, and label-shuffled disease assignments --- probes how
much of that discrimination is real rather than an artifact of the setup, finding the quality-control gate
behaves exactly as designed while the disease hop's stringency is largely inherited from the underlying
enrichment data (Section 4.1b). The
highest-ranked near-novel survivor --- near-novel by an operational criterion of low direct-literature
co-mention plus low curated association, not strict Swanson-style A--C absence --- is \textbf{NAB2}, the worked
example: the perturbation (RNA-seq) data support it as a Th1/Th2 regulator, while its atopic-eczema link is
a \emph{genetic-association nomination} (a GWAS-based disease label, not an expression claim). We further
falsify the sharpest artifactual explanation of NAB2's \emph{perturbation} signal --- a CRISPRi cis-effect on the
adjacent gene \emph{STAT6} --- against the study authors' own genome-wide data; this strengthens the case that
the signal is NAB2-specific, and does not claim to prove the disease link. Wayfinder does not replace
experimental follow-up; it makes the prior question --- \emph{which of these thousands is worth a bench's time?}
--- answerable, with a receipt.

The remainder of this paper is organized as follows. Section 2 provides background on the LBD follow-up
problem, the Perturb-seq substrate, and agentic language models for scientific work. Section 3 details the
generator, the referee, and the workbench. Section 4 reports the verdict ledger, the confident-no and
quality-control behavior, the NAB2 worked example and its confounder falsification, the native
reproduction and self-audit, and an independent cross-model replication. Section 5 discusses what the
agentic loop changes, the next experiment it nominates, and the study's limitations.

\section*{2. Background}\subsection*{2.1 The follow-up problem in literature-based discovery}\label{the-follow-up-problem-in-literature-based-discovery}

Literature-based discovery formalizes a simple observation: two facts can each be established in the
literature while the conclusion that joins them has never been written down. Swanson's original
formulation connects an entity \textbf{A} to an intermediate concept \textbf{B}, and \textbf{B} to an outcome \textbf{C},
where the \textbf{A--C} link is absent \cite{swanson1986}. In \emph{open} discovery one fixes \textbf{A} and searches for
reachable \textbf{C}; in \emph{closed} discovery one fixes both \textbf{A} and \textbf{C} and asks whether a plausible \textbf{B}
bridges them. Four decades of work have refined how the intermediate terms are linked, filtered, and
ranked, and have applied the paradigm across biomedicine --- including our own use of an ABC pipeline to
connect a plasma-metabolite signal to a druggable target in cardiac arrest \cite{henry2021}.

The paradigm's persistent weakness is not generation but \emph{triage}. An ABC search over a large corpus
returns thousands of candidate links, and the ranking signals available to classical LBD are properties
of the literature --- how rare a co-mention is, how strong an intermediate association looks. Rarity in the
literature, however, tracks \emph{obscurity} at least as strongly as it tracks \emph{importance}: a gene can look
novel simply because no one has studied it. The result is a long ranked list that no laboratory can work
through, and the field has said so of itself. Reviewing a proposal submitted for funding consideration,
three reviewers named the problem directly --- LBD \emph{``generates an enormous number of hypotheses, almost none
of which ever get followed up''}. Standard evaluations of LBD (time-slicing, link prediction) measure
whether the generator would have \emph{rediscovered} known links, not whether a specific proposed hypothesis is
supported by independent experimental data it has not seen \cite{henry2021}. We are precise about the scope of the answer this enables: adjudicating a
hypothesis against a held experimental substrate closes the \emph{triage} loop --- which of the machine's
hypotheses is worth pursuing --- not the experimental-follow-up loop itself. What has been missing is that
adjudication step: a way to put each machine-generated hypothesis to a data test and record a verdict ---
including a \emph{no}.

\subsection*{2.2 A held experimental substrate, and what counts as a receipt at each hop}\label{a-held-experimental-substrate-and-what-counts-as-a-receipt-at-each-hop}

The adjudication step needs data that exists independently of the hypotheses being posed. We use a
genome-scale CRISPRi Perturb-seq resource in primary human CD4+ T cells \cite{zhu2025},
in which thousands of genes are individually knocked down and single-cell RNA sequencing reads out the
transcriptional consequence of each perturbation across stimulation conditions. Because the measurements
were made before --- and without reference to --- any triple we test, they function as a \emph{held substrate}: the
referee queries them retrospectively rather than commissioning new experiments. The analysis is
CPU-feasible because it runs over the study's aggregated supplementary tables --- per-guide knockdown
efficiency, differential-expression statistics, a Th1/Th2 polarization signature, and a
cluster-to-autoimmune-disease enrichment --- rather than the raw multi-million-cell matrices. This substrate
is not a universal oracle of biological truth; it is a bounded, independent adjudication surface for
mechanistic T-cell claims --- valuable precisely because, for a hypothesis it \emph{can} speak to, it returns
support, refutation, or \emph{untested} with a receipt, and stays silent (rather than guessing) on the rest.

Crucially, the four tables do not all supply the same \emph{kind} of evidence, and we keep the distinction
explicit throughout. The first three hops --- did the knockdown work, did the perturbation produce a real
transcriptional effect, did that effect shift the T-cell program --- are backed by \textbf{experimental
receipts}: direct measurements from the perturbation data. The fourth hop --- does the perturbed gene's
downstream signature enrich for a disease --- is an \textbf{association receipt}: the disease labels derive from
genetic evidence (GWAS and related association data, without colocalization or linkage-disequilibrium
control), so a positive verdict at this hop is a \emph{nomination}, not a claim of experimental disease
causality. Reading the substrate as a direct experimental readout for the first three hops but an
association signal for the fourth is what keeps the eventual claims calibrated.

\subsection*{2.3 Agentic language models as a scientific instrument}\label{agentic-language-models-as-a-scientific-instrument}

The generation, adjudication, and provenance in this work were carried out inside an \emph{agentic} language-
model environment, and because that setting is newer to the discovery-informatics reader than the biology
or the LBD, we describe what it does and --- more importantly --- what discipline makes it trustworthy for
science. Contemporary language models can be given \emph{tools}: they can call code, query a web API, or read a
file, and act on the returned values rather than on their own recollection. They can also be composed into
an \emph{actor--critic} configuration, in which one model performs an analysis and a second, independent model
reviews it.

Two design commitments make this usable as an instrument rather than a source of confident narration.
First, we separate deterministic \emph{lookups} from \emph{judgment}. Every data receipt in Wayfinder --- a knockdown
efficiency, an effect size, an enrichment odds ratio --- is produced by deterministic code, reproducible and
free of model discretion; the language model \emph{interprets} receipts and assembles the verdict, but it never
asserts a biological fact that a table value does not support. Visible agency is reserved for judgment, not
for data retrieval. Second, the work runs on a scientific workbench --- Claude Science, a persistent-kernel
environment --- in which an author model and an independent reviewer model operate at separate checkpoints:
the reviewer verifies each reported number against the underlying artifacts and enforces calibrated
language on the output, and in this work it flagged and removed overstated words (for example ``validated''
and ``definitive'') from the platform's own text. We are precise about what this independence is and is not:
it is \emph{role, model, and checkpoint} independence --- a distinct reviewer model at distinct verification
points --- within a single model family, not cross-vendor independence. Independence across model families
is provided separately and externally (Section 4.6), by replicating the finding with a different vendor's
models. With that boundary stated, the agentic workbench lets a single researcher run generation,
data-grounded adjudication, and a self-audit of the result's language on one bench --- the capability the
rest of this paper puts to work.

\section*{3. Materials and Methods}\subsection*{3.1 Materials}\label{materials}

\textbf{Perturbation data.} All prioritization runs against the aggregated supplementary tables of a
genome-scale CRISPRi Perturb-seq study (a 2025 preprint) in primary human CD4+ T cells \cite{zhu2025}:
(T1) per-perturbation differential-expression statistics, including an on-target effect flag and a count
of downstream differentially expressed genes; (T2) a Th1/Th2 polarization-signature enrichment per
perturbation, reported for two independent marker contrasts; (T3) a cluster-to-autoimmune-disease
enrichment (odds ratio and FDR per module--disease pair); and (T4) per-guide knockdown efficiency with a
significance call. We work at the 8-hour stimulation condition unless noted. No raw single-cell matrices
are used; for the one confounder check that needs genome-wide differential expression, the study's
deposited matrix is read lazily from public object storage (byte-range reads, no bulk download).

\textbf{Disease and program vocabularies.} The program \textbf{B} is fixed to Th1/Th2 polarization, with two term
arms: process terms (e.g.~``Th2 cells'', ``T helper 2'') used for the gene--program signal, and marker terms
(GATA3, IL-4, IFN-$\gamma$, T-bet) additionally used for the program--disease signal, since disease literature
legitimately cites markers. The disease set \textbf{C} comprises 12 immune diseases, each pinned to a MONDO
identifier \cite{mondo2022} and re-verified live against Open Targets search and the EBI Ontology Lookup Service. This
verification is load-bearing: MONDO and EFO identifiers differ, and a wrong identifier makes Open Targets
report ``no known association'' for every gene, spuriously inflating apparent novelty everywhere.

\textbf{Literature and association sources.} Co-mention counts come from Europe PMC \cite{europepmc2015} \texttt{hitCount} queries (all
terms quoted); curated gene$\rightarrow$disease association scores come from the Open Targets GraphQL API \cite{opentargets2023} --- the
\emph{overall} target--disease association score (not a genetics-only subset), used as a known-link / novelty
prior. This is distinct from the disease \emph{labels} on the enrichment modules (the disease hop, \S{}3.3), which
the source study derives from genetic evidence.

\textbf{Model cast.} Inside the workbench (Section 3.4) the author role is an Opus-class model (Claude Opus
4.8); the independent reviewer role is a Sonnet-class model (Claude Sonnet 5), invoked at separate
verification checkpoints; lightweight inline sampling is handled by a smaller model. All
biological interpretation is confined to these models reading deterministic receipts; no model computes a
receipt.

\subsection*{3.2 Candidate generation}\label{candidate-generation}

\textbf{A disease-answer-free gene universe.} The candidate gene set \textbf{A} is defined from the perturbation
data alone, before any disease information is consulted. A gene enters \textbf{A} if it clears the knockdown-QC
gate ($\geq$1 guide with a significant knockdown call in T4), has a significant transcriptional effect (an
on-target significance flag or a nonzero downstream-DE count in T1), and shows a significant Th1/Th2 shift
(T2 adjusted \emph{p} \textless{} 0.05). Construction never reads the disease table T3, so the gene list cannot be
contaminated by the \emph{disease} answer it will later be tested against; each gene carries its downstream-DE
count forward as an effect-strength prior. \textbf{A} is disease-answer-free, not evidence-free: it is
preselected on the same knockdown/effect/program evidence the referee re-reads at hops 0--2 (a consequence
we make explicit in \S{}3.3). At the 8-hour condition, \textbf{A} contains 3,935 genes.

\textbf{Five deterministic signals.} For each (gene A, disease C) pair we compute: \texttt{ab}, the Europe PMC
co-mention count of A with the Th1/Th2 process terms; \texttt{bc}, the co-mention count of the program terms with
C (constant per disease); \texttt{ac\_known}, the Open Targets \emph{overall} target--disease association score of A for
C (zero when absent --- the novelty signal; the overall score, not a genetics-only subset); \texttt{ac\_lit}, the
direct A--C co-mention count (computed only for referee
survivors, to bound API cost); and \texttt{effect}, the downstream-DE count from T1. Every signal is a
deterministic lookup.

\textbf{Gate and ranking objective.} Eligibility requires all of: \texttt{ab\ $\geq$\ ab\_gate} (a universe percentile,
default the median, which removes candidates that look novel only because the gene is understudied),
\texttt{bc\ $\geq$\ 3}, and \texttt{ac\_known\ $\leq$\ 0.1} (no established curated association). Raw \texttt{ac\_lit} is deliberately \emph{not} in the
gate --- hard-gating on zero co-mention would exclude every well-studied strong regulator and reward
obscurity. Eligible candidates are ranked by a balanced objective,

\begin{equation*}
\mathrm{score} = \min\!\big(z(\mathrm{log1p}\,ab),\; z(\mathrm{log1p}\,bc)\big)
  + \beta\, z(\mathrm{log1p}\,\mathrm{effect})
  - w\,\mathrm{log1p}(\mathit{ac\_lit})
  - w_2\, \mathit{ac\_known},
\end{equation*}

with $\beta$ = 1, w = 1, w2 = 3, and z-scores taken over the universe. The \texttt{min(z\_ab,\ z\_bc)} term makes the
bridge \emph{balanced} --- one strong axis cannot rescue a weak other; \texttt{$\beta$$\cdot$z(effect)} is the ``loud in the data''
reward; the final two terms push novel-yet-bridged candidates up without hard-excluding them. The
objective's structure and weights were fixed before the sweep and were not tuned to any survivor; the
design of this objective is the only human judgment in an otherwise mechanical pipeline. Ranking sets
\emph{priority} among eligible candidates and is separate from the referee's supported/refuted verdict --- a
survivor's verdict does not depend on the weights --- so the choice of weights affects which candidate is
foregrounded, not whether it is supported (we examine rank stability under alternative weights in Section 4).

\textbf{Order of operations.} The sweep builds \textbf{A}, prefetches the per-disease \texttt{bc}/\texttt{ac\_known} (12 calls
each) and per-gene \texttt{ab}, applies the gate, then runs the deterministic referee \emph{first} (it is local and
free) to keep only chain-held survivors (receipt-complete at every hop), and only then computes \texttt{ac\_lit}
and the final score for those survivors. At the 8-hour condition this funnels 3,935 genes $\rightarrow$ 22,039 eligible
(gene, disease) pairs $\rightarrow$ 43 with a positive disease-hop enrichment $\rightarrow$ 30 clean full-chain nominations.

\subsection*{3.3 The referee}\label{the-referee}

The referee evaluates one triple at a time as a four-hop chain, each hop reading a specific receipt:

\begin{itemize}
\tightlist
\item
  \textbf{Hop 0 --- knockdown QC.} Did the perturbation actually work (T4)? If no guide reaches a significant
  knockdown, the triple is returned \textbf{untested} --- never as a negative. This is the gate that converts a
  failed experiment into an honest ``we cannot say,'' rather than a false ``no effect.''
\item
  \textbf{Hop 1 --- effect.} Did the (real) knockdown produce a reproducible transcriptional effect (T1): a
  nonzero downstream-DE count, on target, without an off-target flag?
\item
  \textbf{Hop 2 --- program.} Do the downstream effects shift the Th1/Th2 program (T2)?
\item
  \textbf{Hop 3 --- disease.} Does the perturbed gene's downstream signature fall in a module enriched for the
  specific disease (T3: odds ratio, FDR)? This is an \textbf{association / enrichment} receipt --- the module's
  disease label is genetic (GWAS-based) --- not experimental evidence of disease causality.
\end{itemize}

A triple is \textbf{supported} only if the full chain holds with a nonzero, positive effect; effect-zero chains
are demoted to \textbf{supported-weak}, off-target-flagged chains to \textbf{supported-flagged}, a failed effect hop
yields \textbf{refuted-effect}, and failure at the specific-disease hop is \textbf{refuted-for-C}. Throughout,
\textbf{supported} is a bounded \emph{verdict label} --- the chain held with a receipt at each hop, the disease hop
remaining a genetic-association nomination --- not a claim that the biology is proven. The confident \emph{no}
lives in these demotions, all deterministically computed.

We separate two categories that the rest of the paper --- and Figure 1 --- keep visually and textually
distinct: \textbf{construction filters} (the knockdown-QC, effect, and program pre-gates that build the
universe \textbf{A}) and \textbf{referee decisions} (the disease-C specificity test, together with the QC and
effect demotions the referee adjudicates on \emph{arbitrary} genes). The referee genuinely owns the QC and
effect hops; the disease hop it can only \emph{prioritize} as a genetic-association nomination. Two honesty
points follow from this separation, and belong in Methods, not only in Limitations. First, because \textbf{A}
is preselected on knockdown-QC, effect, and program-significance, \emph{within the funnel} hops 0--2 are
largely pre-gated (the
program hop is a strict tautology --- no eligible gene can fail it), so among funnel survivors the dominant
cull is disease-C specificity rather than broad four-hop falsification. The confident-\emph{no} behaviour is
therefore shown two ways: by the referee's demotions and exact-disease refutations within the funnel, and
--- more tellingly --- by applying the same referee to \emph{arbitrary} genes outside the pre-gated universe, where
it discriminates at every hop (a gene whose knockdown fails QC returns \emph{untested} at hop 0; a gene with no
real transcriptional effect is \emph{refuted} at hop 1). That out-of-funnel ledger is reported in Section 4.
Second, the program hop, being a within-funnel tautology, discriminates in an individual triple's receipt
--- as a reported quantity --- not as an independent filter on funnel survivors.

\subsection*{3.4 The agentic loop, and how the instrument is operated}\label{the-agentic-loop-and-how-the-instrument-is-operated}

Generation, evidence integration, and provenance assembly were run inside Claude Science, a persistent-kernel
agentic scientific workbench, in which an author model writes and runs analysis code and an independent
reviewer model verifies the result at separate checkpoints (Section 3.1).

\textbf{Operating an instrument that has no API.} Claude Science exposes no programmatic task-submission
interface: analyses are ordinarily driven by hand in its web interface, which makes each run a one-off,
hand-paced event rather than a reproducible procedure --- human clicking becomes the bottleneck at exactly
the moment when keeping pace with the twin floods of data and literature is the whole point. To make the
workbench \emph{scriptable}, we drive its web interface
headlessly. A browser-automation driver, orchestrated from a coding agent (Claude Code), loads an operator-owned
authenticated session, submits the task prompt, and \textbf{auto-approves the workbench's in-loop sandbox
prompts --- working-directory, code-execution, and network-access cards --- for unattended, zero-click
operation}; it then polls the run to completion and retrieves the emitted artifacts together with the
workbench's own audit records. The auto-approval replaces a human \emph{click}, not a safety boundary: the run
executes in the workbench's own sandbox, in a fixed workspace, under an operator-owned session, and every
approval is logged in the audit store --- so the consent trail is inspectable after the fact rather than
absent. Network approvals were granted only for the named data and literature endpoints the task used --- an
auditable, endpoint-specific approval trail rather than a claim of an enforced allowlist --- and, as the
liveness accounting below makes explicit --- the full-scale run reads from a local cache under a guard that
raises on any live call, so unattended approval cannot silently alter the data receipts. We are candid that
this is browser automation against a user interface with no stable public contract (Section 5.3); it is
nonetheless what converts a click-once web session into a re-runnable, audited pipeline, and --- where a
platform exposes no programmatic interface --- an under-appreciated route to replicable-in-principle
(UI-dependent) \emph{agentic} science.

\textbf{Provenance and self-audit.} Every run is captured in the workbench's own audit store, from which we
draw model identities, per-step costs, and the reviewer's checks. The reviewer verifies each reported
number against the underlying artifacts and enforces calibrated language on the manuscript-facing output ---
the receipt chain and the verdicts a reader sees; in this work it flagged overstated words (``validated'',
``definitive'') in that output and they were removed. This self-audit is \emph{language hygiene plus
receipt-consistency} --- that the wording stays calibrated and each cited number matches its artifact --- not
an independent epistemic verification that the underlying biology is correct; that check is the held
substrate's to make, not the critic's. We scope this claim honestly: the enforcement applies
to the manuscript-facing receipt chain, not to every intermediate machine-generated log (an earlier
generation record still carries a legacy ``validated'' string in a free-text field), which is why
manuscript-facing verdicts cite the post-review receipt rather than the raw generation log. We characterize
the independence precisely, too: it is \emph{role, model, and checkpoint} independence --- a distinct reviewer
model at distinct verification points --- within a single model family, not independence across vendors.

\textbf{Three claims about liveness, kept distinct.} We separate what ran live from what was replayed, because
conflating them would overstate the result. (i) \emph{Full-scale reproduction}: the generation pipeline ran
unchanged over the entire 3,935-gene universe inside the workbench, but against a workbench-local cache of
previously captured literature/association responses, under a guard that raises on any live network call
(cache delta zero) --- a faithful recomputation, not a live crawl. (ii) \emph{Live authorship}: in a separate,
smaller run the workbench was given only the method and wrote its own generator from scratch, making live
Europe PMC and Open Targets calls; the headline gene did not appear in that small run's top candidates, so
nothing was cherry-picked, and a Stage-0 probe independently confirmed live external access from the
kernel. (iii) \emph{Cross-family independence}: replication with a different vendor's models is kept external
by design (Section 4.6). Reported this way, the workbench supplies generation, data-grounded
prioritization, and a language self-audit; cross-model independence is supplied from outside it.

\section*{4. Results}We report six results. The funnel and verdict ledger show the referee discriminating across a spread of
worked examples (\S{}4.1), and two negative controls plus a rank-stability check probe how much of that
discrimination is a property of the data versus an artifact of the setup --- with one control returning a
deliberately honest, non-obvious answer (\S{}4.1b) --- and a gate-threshold grid shows the flagship's rank is
stable wherever it stays eligible (\S{}4.1c). The confident \emph{no} and the quality-control
catch are examined on their own (\S{}4.2), because they are the point. NAB2 is then followed hop by hop as the
worked hero (\S{}4.3), and its sharpest artifactual confounder is falsified against the study authors' own
genome-wide data (\S{}4.4), while the one confounder the substrate \emph{cannot} discharge --- the 12q13
GWAS-label provenance --- is foregrounded rather than buried (\S{}4.4b). Finally we show the entire loop reproduced natively inside the agentic workbench,
audited by an independent critic model (\S{}4.5), and replicated by a second vendor's models (\S{}4.6).

\subsection*{4.1 The funnel and the verdict ledger}\label{the-funnel-and-the-verdict-ledger}

At the 8-hour stimulation condition, the generator built a disease-answer-free universe of \textbf{3,935}
knockdown-gated, program-significant genes, from which the gate admitted \textbf{22,039} eligible
(gene, disease) pairs. Running the deterministic referee over those pairs returned \textbf{43} with a positive
disease hop, of which \textbf{30} held the full chain cleanly (nonzero positive effect); the remainder split into
10 effect-zero \emph{supported-weak}, 3 off-target \emph{supported-flagged}, and 1 \emph{refuted-effect}, with the other
\textbf{21,995} pairs \emph{refuted for the specific disease} (Table 3). The generate-to-cull ratio is the headline
shape: 22,039 machine-generated questions reduced by data receipts to 30 clean survivors, none asserting
more than a table value supports.

We are deliberately careful about what this funnel does and does not demonstrate, and we restate the two
honesty caveats here rather than defer them to Limitations. First, the universe \textbf{A} is preselected on
program-significance, so \emph{within the funnel} the program hop cannot fail (\texttt{refuted\_program\ $\equiv$\ 0} by
construction); the dominant cull among eligible pairs is therefore disease-specificity, not broad four-hop
falsification. The program hop discriminates in an individual triple's receipt (a reported quantity --- NAB2's
Ota \emph{z} = 7.71), not as an independent filter on funnel survivors. Second, the ``43'' is a \textbf{joint} product
of the literature-novelty gate and the referee, not a referee-only count: the referee alone supports
\textbf{395 of 47,220} pairs, and the novelty gate culls 395 $\rightarrow$ 43. The honest reading is ``referee-supported
\emph{among literature-eligible} pairs.'' Accordingly we present the ledger below as a \textbf{demonstration that the
referee discriminates} --- that on worked examples it will support, decline to test, and refute, each with a
receipt --- not as a measured precision or recall.

The verdict ledger (Table 2) places five genes across the full verdict spectrum, each with its per-hop
receipt:

{\def\LTcaptype{none} % do not increment counter
\begin{longtable}[]{@{}
  >{\raggedright\arraybackslash}p{(\linewidth - 4\tabcolsep) * \real{0.3333}}
  >{\raggedright\arraybackslash}p{(\linewidth - 4\tabcolsep) * \real{0.3333}}
  >{\raggedright\arraybackslash}p{(\linewidth - 4\tabcolsep) * \real{0.3333}}@{}}
\toprule\noalign{}
\begin{minipage}[b]{\linewidth}\raggedright
Gene $\rightarrow$ disease (condition)
\end{minipage} & \begin{minipage}[b]{\linewidth}\raggedright
Verdict
\end{minipage} & \begin{minipage}[b]{\linewidth}\raggedright
Deciding receipt
\end{minipage} \\
\midrule\noalign{}
\endhead
\bottomrule\noalign{}
\endlastfoot
\textbf{NAB2} $\rightarrow$ atopic eczema (Stim8hr) & \textbf{supported} & full chain holds; eczema clusters OR 3.90 / FDR 0.0028 and OR 3.43 / FDR 0.0224 (\S{}4.3) \\
\textbf{EGR2} $\rightarrow$ asthma (Stim8hr) & \textbf{supported} & 2/2 guides signif; effect -11.06, 854 downstream DE; member of 34 disease clusters, asthma OR 20.4 / FDR 5$\times$10\textsuperscript{-}\textsuperscript{5} \\
\textbf{NUDT1} $\rightarrow$ type 1 diabetes (Stim8hr) & \textbf{supported-weak} & full chain holds but a trivial effect (4 downstream DE) --- the \emph{only} pure-disjoint (zero-literature) survivor \\
\textbf{IL2} $\rightarrow$ (Rest) & \textbf{untested} & knockdown-QC fails: 0/2 guides significant; target barely expressed at Rest, nothing to knock down \\
\textbf{SLC1A5} $\rightarrow$ type 1 diabetes (Stim8hr) & \textbf{refuted} & in 9 disease cluster gene-sets, but none reach FDR \textless{} 0.05 (best FDR 0.054, OR 2.73) \\
\end{longtable}
}

Two features of the ledger carry the thesis. The \emph{untested} verdict (IL2) is not a missing entry or a
failure to compute --- it is the referee refusing to read an uninterpretable knockdown as a biological
result. And the \emph{refuted} verdict (SLC1A5) is a plausible metabolite-transporter $\rightarrow$ autoimmunity link that
the enrichment data decline to support: a confident, receipt-backed \emph{no}. NUDT1 illustrates the opposite
edge --- a hypothesis that clears every hop yet is foregrounded only weakly, because its effect is trivial;
this is exactly why the pipeline \emph{ranks} by a novelty-plus-effect score rather than hard-gating on
zero-literature absence, which would reward obscurity (\S{}3.2). We note that these examples are drawn to span
the verdict space, not sampled to estimate a rate; \S{}4.1b supplies a disease-hop support-\emph{rate} control over
the full A$\times$C space --- not a ledger accuracy rate, nor a claim of precision or recall.

\subsection*{4.1b Sensitivity: two negative controls and rank stability}\label{b-sensitivity-two-negative-controls-and-rank-stability}

To move ``the referee discriminates'' from assertion to a small measured result, we ran three checks as a
bounded kernel analysis inside the workbench (\S{}3.4).

\textbf{Control 1 --- the QC gate does its job.} We applied the referee to every gene whose knockdown \emph{fails} QC
at the 8-hour condition (no guide reaching a significant knockdown call in T4): \textbf{2,430} genes. All 2,430
returned \emph{untested} at hop 0 --- a \textbf{100\%} rate, with no leakage into a spurious \emph{refuted} or \emph{supported}
verdict. The gate converts failed experiments into an honest ``cannot say,'' exactly as designed, and does so
without exception across the full failed-knockdown set.

\textbf{Control 2 --- the disease hop's stringency is substrate-inherited, and its pass count is label-dependent.}
Across the referee's native pair space of all \textbf{47,220} (gene $\in$ A) $\times$ (disease $\in$ C) combinations at the
8-hour condition, the disease hop supports only \textbf{406 (0.86\%)} and refutes the remaining \textbf{99.14\%}. That
stringency, however, is not the referee's own doing: under \textbf{2,000} permutations of the disease labels
across the enrichment rows, the null still passes only \textbf{467.7 $\pm$ 10.9} pairs --- a null pass rate of
\textbf{0.99\%} --- so even randomly-labelled pairs clear FDR \textless{} 0.05 barely more than one percent of the time. The
near-total refutation is therefore \textbf{inherited from the enrichment study's FDR family} --- a
substrate-provenance property we flag in Limitations (\S{}5.3), and one an adversarial replication
independently noted --- not a discrimination we should credit to the referee. What the permutation \emph{does}
establish is that the pass count is \textbf{label-dependent}: the observed 406 sits \textbf{5.6 standard deviations
below} the null mean (lower-tail permutation \emph{p} $\approx$ 5$\times$10\textsuperscript{-}\textsuperscript{4}; the \emph{upper}-tail \emph{p} is 1.0 --- that is, the hop
does \textbf{not} pass \emph{more} often than chance, so we explicitly do not claim a ``rarer-than-chance''
selectivity). The direction --- observed \emph{below} null --- is a feature of this restricted twelve-disease
permutation setup; the one component we actually measured, a per-gene concentration effect, is small (a
passing gene carries \textbf{2.39} distinct disease labels under the true assignment versus \textbf{2.44} under
shuffling), so we report the figures but do not interpret the mechanism of the direction further. The load-bearing reading is the
one we can defend on both tails: the disease hop is a \textbf{stringent, label-dependent nomination filter whose
stringency is substrate-inherited} --- not a demonstration of the referee's own discriminating power. (This
control speaks to the disease hop; it does not re-open the program-hop tautology of \S{}3.3, which is a
construction property, not a chance one.)

\textbf{Rank stability under alternative objective weights.} The ranking objective (\S{}3.2) has three free weights
($\beta$ on effect, \emph{w} on direct co-mention, \emph{w\textsubscript{2}} on curated association; default $\beta$ = 1, \emph{w} = 1, \emph{w\textsubscript{2}} = 3). The
weights set \emph{priority} among eligible survivors, not the supported/refuted verdict, which is
weight-independent by construction --- but a headline example that is rank-4 only at one weight setting would
be a fragile choice of hero. We recomputed the ranking over a \textbf{27-point grid} crossing $\beta$ $\in$ \{0.5, 1, 2\},
\emph{w} $\in$ \{0.5, 1, 2\}, and \emph{w\textsubscript{2}} $\in$ \{1, 3, 5\}, recovering each survivor's weight-independent score component from
its default-weight value and re-ranking under every setting (the default setting reproduces NAB2's rank of
4, confirming the recovery). NAB2's rank ranged from \textbf{1 to 8} across the grid, with a \textbf{median of 4}, and
it stayed within the top five survivors under \textbf{24 of the 27} settings (89\%). No survivor's \emph{verdict}
changed, since verdicts do not depend on the weights. The worked hero is robust to the one human judgment
(\S{}3.2) in an otherwise mechanical pipeline.

\subsection*{4.1c Threshold robustness: survivors and the flagship rank across the gate}\label{threshold-robustness-survivors-and-the-flagship-rank-across-the-gate}

Section 4.1b varied the ranking \emph{weights}; the eligibility \emph{gate} (\S{}3.2) has its own three
thresholds --- the literature-percentile floor \texttt{ab\_gate\_pct} (default the median, 0.50), the
program--disease co-mention floor \texttt{min\_bc} (default 3), and the novelty ceiling \texttt{tau} on
curated association (default 0.10). We swept a \textbf{27-point grid} crossing \texttt{ab\_gate\_pct} $\in$
\{0.25, 0.50, 0.75\}, \texttt{min\_bc} $\in$ \{2, 3, 5\}, and \texttt{tau} $\in$ \{0.05, 0.10, 0.20\},
recomputing the funnel offline for every cell. Because a pair's referee verdict does not read the gate, the
verdicts are computed once over the full 47,220-pair space and re-filtered per cell (the default cell
reproduces the 30 clean survivors and NAB2's rank of 4 exactly). We read three things from the grid, as a
diagnostic of the machinery's response to its thresholds, not as a validation.

\emph{The survivor count behaves as a gate should.} The clean full-chain survivor count ranges from
\textbf{7} at the strictest cell to \textbf{74} at the loosest, rising as the gate is relaxed, with the
default setting (30) in mid-range --- the gate admits more candidates as it loosens, without discontinuity.

\emph{The flagship's rank is stable wherever it is eligible.} Whenever NAB2 $\rightarrow$ atopic eczema clears
the gate, it ranks in a tight band --- \textbf{1, 4, or 5} across every such cell, median \textbf{4}, never
worse than fifth --- echoing the rank-4 median of the weight sweep (\S{}4.1b). Its \emph{verdict} is
invariant: the referee supports the pair in all 27 cells.

\emph{But eligibility depends on the literature floor, and we say so plainly.} NAB2 survives as a clean
nomination in \textbf{18 of 27} cells, and the nine exceptions are exactly the nine at
\texttt{ab\_gate\_pct}~=~0.75 --- independent of \texttt{min\_bc} and \texttt{tau}. NAB2's A--B program-literature
co-mention sits between the median and the 75th percentile of the universe, so a floor at the 75th percentile
removes it from the \emph{candidate pool} --- not because its receipts fail (they do not) but because a
near-novel gene, by construction, carries only moderate program-literature. This is the honest cost of
ranking near-novelty: the flagship is robust to the two disease-side thresholds (\texttt{min\_bc},
\texttt{tau}) but sensitive to how aggressively the literature floor prunes understudied genes, a dependence
a reader should weigh alongside the nomination. The clean-survivor \emph{sets} overlap the default by a
Jaccard of \textbf{0.41} at the median cell (minimum 0.23) --- a gate that reshapes the candidate pool around
a stable core. (This grid is deterministic and offline; recomputing the survivor sets and NAB2's rank reads
only cached literature counts.)

\subsection*{\texorpdfstring{4.2 The confident \emph{no}, and the quality-control catch}{4.2 The confident no, and the quality-control catch}}\label{the-confident-no-and-the-quality-control-catch}

The moat of this approach is a referee that will answer \emph{no} --- and that distinguishes a \emph{refuted}
hypothesis from an \emph{untested} one. Both behaviors appear in the ledger, and both are worth stating plainly.

IL2 is the quality-control catch. IL2 is a canonical, biologically central T-cell gene that appears in both
the Th program and the disease gene-sets; a naive pipeline that read the empty downstream result of a
failed knockdown would call it ``no effect'' and record a false negative. The referee instead halts at hop 0:
0 of 2 guides reach a significant knockdown (guide mean expression 0.031 vs.~non-targeting-control 0.036 ---
the target is barely expressed at Rest, so there is nothing to knock down), and the verdict is \emph{untested ---
knockdown failed QC}. An artifact is caught, not mistaken for biology.

SLC1A5 is the correctly-sourced refutation. A metabolite transporter with a plausible metabolic-gene $\rightarrow$
autoimmunity story, SLC1A5 clears the QC gate (1/2 guides significant) and shifts the Th1/Th2 program, yet
its disease hop \emph{refutes}: it appears in nine autoimmune-disease cluster gene-sets but none reaches
FDR \textless{} 0.05 (best FDR 0.054, type 1 diabetes, OR 2.73). The enrichment data decline the link. This is the
verdict a confirmation-shaped pipeline never produces, and it is deterministic --- the \emph{no} is computed, not
argued.

\subsection*{4.2b Out-of-funnel discrimination: the referee's own edge, quantified}\label{out-of-funnel-discrimination-the-referees-own-edge-quantified}

Section 4.1b established that the disease hop's near-total refutation rate is \emph{substrate-inherited} --- a
property of the enrichment study's FDR family, not the referee's own doing. That invites the fair question a
referee report pressed: once the substrate-inherited disease hop is set aside, how much does the referee
\emph{itself} discriminate? Its own hops are the knockdown-QC gate (hop 0) and the effect and program checks
(hops 1--2), each a deterministic computation the referee owns. We quantify their discrimination two ways,
both offline and outcome-independent.

\emph{On arbitrary genes.} We applied the referee to \textbf{all 11,415} genes perturbed at the 8-hour
condition --- the full out-of-funnel population, not the pre-gated universe --- pairing each with an arbitrary
disease (the referee's hops 0--2 are disease-independent). The referee's own hops cull \textbf{16.9\%} of
them: \textbf{16.8\%} (1{,}914 genes) return \emph{untested} because the knockdown fails QC --- genes the
substrate cannot speak to --- with only a small remainder refuted at the effect or program hop. The
knockdown-QC gate is thus the referee's dominant own-edge discriminator (one arbitrary gene in six cannot be
tested here), while the effect gate, lenient by design (any nonzero downstream signal passes), rarely fires
on unselected genes. This is the honest shape of the referee's own \emph{no}: it is real, and it is
concentrated in the quality-control catch rather than spread across the chain.

\emph{On a frozen nominator's best guesses (hard negatives).} A stronger test must not choose genes by the
outcome being evaluated. We froze a naive curated-association nominator --- for each of the 12 diseases, the
top 50 perturbed genes by Open Targets association score, selected with no reference to the referee --- and
applied the referee to the resulting \textbf{600} association-plausible (gene, disease) nominations. The
verdicts decompose exactly along the line the reviews asked us to draw: \textbf{15.7\%} are culled at the
referee's \emph{own} hops (\textbf{75} \emph{untested}, 8 \emph{refuted-effect}, 11 \emph{refuted-program}),
\textbf{67.3\%} are \emph{refuted for the specific disease} (the substrate-inherited hop, \S{}4.1b), and only
\textbf{17.0\%} are supported. The own-hop culls are the hard negatives --- genes a database-driven nominator
would confidently advance that the data cannot support here --- and they include strongly-associated flagship
pairs: \emph{IL36RN} $\times$ psoriasis (association 0.82), \emph{TREX1} $\times$ systemic lupus (0.78), and
\emph{PADI4} $\times$ rheumatoid arthritis (0.68), each returned \emph{untested} because its knockdown failed
QC in this screen. A curated-association ranking, however strong, cannot tell you whether the perturbation is
even measurable; the referee's own hops can, and they flag roughly one in six of the top nominations. This is
the discrimination the referee supplies on its own --- reported separately from, and not conflated with, the
disease hop's substrate-inherited stringency.

\subsection*{4.3 NAB2 $\rightarrow$ Th1/Th2 $\rightarrow$ atopic eczema, hop by hop}\label{nab2-th1th2-atopic-eczema-hop-by-hop}

The highest-ranked near-novel survivor is \textbf{NAB2}, near-novel by an operational criterion --- low direct
literature co-mention (\texttt{ac\_lit} = 6) and no curated Open Targets association (\texttt{ac\_known} = 0.038) --- rather
than strict Swanson-style A--C absence. The data support NAB2 as a Th1/Th2 regulator with a receipt at every
hop (Figure 4):

\begin{itemize}
\tightlist
\item
  \textbf{Hop 0 --- knockdown QC.} 2/2 guides significant, best adjusted \emph{p} $\approx$ 1$\times$10\textsuperscript{-}\textsuperscript{1}\textsuperscript{6} (guide expression 0.056
  vs.~non-targeting control 0.567). The perturbation worked.
\item
  \textbf{Hop 1 --- effect.} On-target knockdown, effect -16.9 (-16.88 exact), \textbf{301 downstream differentially expressed genes},
  no off-target flag. A large, clean transcriptional consequence.
\item
  \textbf{Hop 2 --- program.} Th2-associated in the Ota contrast \cite{ota2021} (\emph{z} = 7.71, adjusted \emph{p} = 1.96$\times$10\textsuperscript{-}\textsuperscript{1}\textsuperscript{3},
  log-fold-change +0.63; polarity marker-checked); not significant in the H\"{o}llbacher contrast \cite{hollbacher2020} --- one of
  two contrasts, reported rather than hidden.
\item
  \textbf{Hop 3 --- disease.} Member of two atopic-eczema-enriched modules: OR 3.90 (FDR 0.0028) and OR 3.43
  (FDR 0.0224). These are genome-wide functional immune modules (significant clusters 90 and 100 ---
  BACH2, BCL6, IRF4, CD28, IL4, IL10), not a 12q13 locus artifact.
\end{itemize}

The precisely scoped claim is: \emph{consistent with a re-derived NAB2 $\rightarrow$ Th1/Th2 $\rightarrow$ atopic-eczema chain the literature
has not made.} We keep the two receipt classes explicit even for the hero --- the knockdown, effect, and
program hops rest on experimental perturbation measurements; the atopic-eczema link is a \textbf{genetic-
association nomination} (an Open Targets GWAS-based disease label, without colocalization or LD control),
not an expression claim or a proof of disease causality. An independent four-agent literature audit found
zero papers connecting NAB2 either to Th1/Th2 polarization or to atopic eczema --- the finding is novel in
the literature --- and the same audit surfaced the confounder we take up next.

\subsection*{\texorpdfstring{4.4 Falsifying the hardest confounder: the STAT6 \emph{cis}-effect}{4.4 Falsifying the hardest confounder: the STAT6 cis-effect}}\label{falsifying-the-hardest-confounder-the-stat6-cis-effect}

NAB2 sits \textasciitilde1.9 kb from \emph{STAT6}, the master type-2/atopic regulator, so the sharpest artifactual explanation
of NAB2's \emph{perturbation} signal is a CRISPRi \emph{cis}-effect: a guide targeting NAB2 could inadvertently
repress the adjacent STAT6, and the observed Th2/eczema signal could be STAT6's, not NAB2's. We test this
directly against the study authors' own deposited genome-wide differential-expression matrix (the 16.8 GB
\texttt{GWCD4i.DE\_stats.h5ad}, read lazily from public object storage by byte-range, no download). Under NAB2
knockdown at Stim8hr:

{\def\LTcaptype{none} % do not increment counter
\begin{longtable}[]{@{}
  >{\raggedright\arraybackslash}p{(\linewidth - 6\tabcolsep) * \real{0.2000}}
  >{\raggedleft\arraybackslash}p{(\linewidth - 6\tabcolsep) * \real{0.2667}}
  >{\raggedleft\arraybackslash}p{(\linewidth - 6\tabcolsep) * \real{0.2667}}
  >{\centering\arraybackslash}p{(\linewidth - 6\tabcolsep) * \real{0.2667}}@{}}
\toprule\noalign{}
\begin{minipage}[b]{\linewidth}\raggedright
Gene (under NAB2-KD)
\end{minipage} & \begin{minipage}[b]{\linewidth}\raggedleft
log\textsubscript{2}FC
\end{minipage} & \begin{minipage}[b]{\linewidth}\raggedleft
adj. \emph{p}
\end{minipage} & \begin{minipage}[b]{\linewidth}\centering
moved?
\end{minipage} \\
\midrule\noalign{}
\endhead
\bottomrule\noalign{}
\endlastfoot
\textbf{STAT6} (neighbor under test) & \textbf{+0.087} & \textbf{0.788} & no --- \textbf{unmoved} \\
\textbf{NAB2} (self / on-target) & \textbf{-3.078} & 7.2$\times$10\textsuperscript{-}\textsuperscript{6}\textsuperscript{0} & yes --- knockdown worked \\
\end{longtable}
}

Of 10,282 measured genes, 302 are significantly moved by NAB2 knockdown ($\approx$ the referee's effect count of
301); STAT6 is not among them, ranking 5,444/10,282 by absolute log-fold-change --- the less-affected half. A
\emph{cis}-artifact would push STAT6 \emph{down}; it does not move. Three further checks corroborate this --- drawn from
the study's own QC fields and our finding analysis rather than the genome-wide DE row above: NAB2 and STAT6
share zero perturbation-effect clusters (their knockdowns are not equivalent); NAB2 clears the study's own
reproducibility bar (cross-guide/donor R 0.74); and the authors' own off-target flag --- a transcription
start site within 10 kb showing significant down-regulation --- is False for NAB2, meaning their pipeline
already found no \emph{cis} down-regulation of any near neighbor, STAT6 included. We therefore report the
\emph{cis}-artifact as \textbf{refuted}: the signal is consistent with NAB2 perturbation rather than STAT6 bleed. We
scope the claim precisely --- this strengthens the case that the \emph{perturbation} signal is NAB2-specific; it
does not prove the disease link, which remains a genetic-association nomination.

\subsection*{\texorpdfstring{4.4b The confounder we cannot discharge: the 12q13 atopy locus}{4.4b The confounder we cannot discharge: the 12q13 atopy locus}}\label{the-confounder-we-cannot-discharge-the-12q13-atopy-locus}

The STAT6 falsification addresses one of \emph{three} distinct ways the 12q13 locus could manufacture NAB2's
eczema signal, and we separate them deliberately because only two are answerable with our data --- and we do
not want the one we can falsify to stand in for the one we cannot. NAB2 lies within the 12q13
atopic-dermatitis susceptibility locus, \textasciitilde1.9 kb from \emph{STAT6}. \textbf{(i) A CRISPRi \emph{cis}-effect}
--- the NAB2 guide silently repressing the adjacent STAT6, so the transcriptional signal is STAT6's --- is
\textbf{falsified}: STAT6 is unmoved under NAB2 knockdown (\S{}4.4). \textbf{(ii) A cluster-membership
artifact} --- NAB2's eczema modules being a 12q13-colocated gene set rather than genome-wide functional
modules --- is \textbf{rejected}: the significant modules (clusters 90 and 100) are genome-wide functional
immune modules (BACH2, BCL6, IRF4, CD28, IL4, IL10), with STAT6 absent (\S{}4.6). \textbf{(iii) An
LD-inherited disease label} we \textbf{cannot} settle. The disease label on those modules is a GWAS-based
genetic association, assigned without colocalization or linkage-disequilibrium analysis; because NAB2 sits
inside the 12q13 atopy locus, its eczema association could in principle be linkage-inherited from the locus's
lead signal --- plausibly STAT6's --- rather than being NAB2-specific. Our substrate carries the aggregated
enrichment result, not the variant-level GWAS, so it \emph{cannot} perform the colocalization that would
separate a genuinely NAB2-specific association from an LD shadow. We therefore \textbf{foreground this as the
flagship nomination's key open question, and explicitly do not claim it discharged.} The perturbation
evidence for NAB2 as a Th1/Th2 regulator is receipt-backed and confounder-tested; the further step from that
regulatory role to \emph{atopic eczema specifically} rests on a genetic-association label whose locus
provenance remains unresolved. Resolving it requires variant-level colocalization against the 12q13
atopic-dermatitis GWAS --- outside this substrate, and precisely the disease-hop analysis the loop nominates
next (\S{}5.2, \S{}5.3). Stating the boundary this plainly is itself the calibrated-language discipline the
paper argues for: the honest verdict here is \emph{a receipt-backed regulatory nomination with an unresolved
disease-label provenance}, not a settled NAB2$\rightarrow$eczema link.

\subsection*{4.5 Native reproduction and self-audit in Claude Science}\label{native-reproduction-and-self-audit-in-claude-science}

The entire loop --- generation, evidence integration, and the assembly of its own provenance --- was run inside the
agentic workbench, at a total run cost of \textasciitilde\$6.41, with the model identities, per-step costs, and reviewer
checks drawn from the workbench's own audit store. We keep the three liveness claims distinct (\S{}3.4). \emph{Full-
scale reproduction}: the real generator ran unchanged over the entire 3,935-gene universe inside the
workbench and reproduced the exact funnel (3,935 $\rightarrow$ 22,039 $\rightarrow$ 43 $\rightarrow$ 30) and NAB2's rank of 4, but against a
workbench-local cache of previously captured literature/association responses under a guard that raises on
any live network call (cache delta zero) --- a faithful recomputation, not a live crawl. \emph{Live authorship}: a
separate 12-gene from-scratch run, given only the method, had the workbench write its own generator and make
live Europe PMC and Open Targets calls; NAB2's effect was too low to rank in that small run's top twelve, so
it did not appear --- nothing was cherry-picked --- and independent re-querying of the workbench's recorded
queries against the live APIs returned identical counts. \emph{Cross-family independence} is kept external by
design (\S{}4.6).

The self-audit is the part that mirrors the paper's own thesis (Figure 5). The author model (an Opus-class
model) assembled the end-to-end receipt chain; an independent reviewer model (a Sonnet-class model),
invoked at separate checkpoints, verified every reported number against the underlying artifacts and
enforced calibrated language on the manuscript-facing output. It produced four verification checks: two
confirming that every cited receipt and the final ``what ran where'' table matched the source of truth, and ---
the telling one --- a check that \textbf{flagged the words ``validated'' (in a title) and ``definitive'' (in a heading)
as calibrated-language violations, both of which were then removed}. The platform corrected its own
overstatement. We characterize this independence honestly: it is \emph{role, model, and checkpoint} independence
--- a distinct reviewer model at distinct verification points within a single model family --- not independence
across vendors. It is also why manuscript-facing verdicts cite the post-review receipt rather than the raw
generation log, which still carries a pre-critic ``validated'' string in a free-text field: that pre- to
post-critic difference \emph{is} the self-audit.

\subsection*{4.6 Independent cross-model replication}\label{independent-cross-model-replication}

Cross-family independence was supplied from outside the workbench, by a five-member replication lab: three
Opus-class agents and two Codex (a different vendor's) agents, each re-deriving the finding from the raw
supplementary tables under an adversarial mandate, with two members performing clean-room re-
implementations that imported none of our pipeline code. The verdict was a \textbf{unanimous pass}. Across the
lab, every headline number was independently reproduced to the unit --- not every member re-checked every
number, but each number was re-derived from the raw CSVs by at least one clean-room member: the NAB2 receipt (2/2 guides,
effect -16.88, 301 downstream DE, Ota \emph{z} +7.71, eczema clusters OR 3.90/3.43), the full funnel
(3,935 $\rightarrow$ 22,039 $\rightarrow$ 43 $\rightarrow$ 30, with the 30/10/3/1 class split and the single pure-disjoint survivor), and the
answer-free construction of the universe.

The replication is worth reporting precisely \emph{because} it was not silent agreement: the adversarial pass
caught and corrected real errors. It found a cluster-ID misalignment in our confounder script --- sorted
cluster IDs printed against row-order FDRs had mislabeled NAB2's significant eczema modules as 74/90 when
the true significant modules are 90 and 100 (cluster 74 is non-significant, FDR 0.52), and the locus test
had been run on the wrong cluster; the corrected run confirms modules 90 and 100 are genome-wide functional
immune modules with STAT6 absent, which rejects the \emph{cluster-membership} locus artifact --- though not the
distinct question of whether the GWAS disease \emph{label} is LD-inherited from the 12q13 atopy locus, which the
substrate cannot settle and which we leave open (\S{}5.3). It corrected an effect-size overstatement (``\textasciitilde8$\times$ stronger'' is the
log-fold-change ratio; the \emph{z}-ratio is \textasciitilde3$\times$). And it reframed the distinctness argument onto its strongest
support --- cluster co-membership, magnitude, and guide-specificity --- rather than the shared disease profile,
which if anything \emph{aids} the confounder. That a hostile independent re-derivation reproduced every number
and improved the framing is the strongest available evidence that the finding is a property of the data,
not of our implementation.

\section*{5. Discussion}\subsection*{5.1 What the agentic loop changes}\label{what-the-agentic-loop-changes}

We opened with a critique the field has leveled at itself --- that literature-based discovery ``generates an
enormous number of hypotheses, almost none of which ever get followed up.'' Wayfinder is a response to that
critique, and it is worth being exact about \emph{which} part of it we answer. We do not claim to have solved
experimental follow-up; the wet-lab confirmation of any surviving hypothesis remains future work, as it
must. What the approach changes is the step immediately before that: \textbf{triage}. Of 22,039
machine-generated gene$\rightarrow$program$\rightarrow$disease hypotheses, the deterministic referee reduced the set to a few
receipt-backed survivors --- and, crucially, attached to each survivor a per-hop receipt and to each
rejection a reason. ``Which of these thousands is worth a bench's time?'' stops being a question answered by
literature-rarity heuristics (which, as we showed, track obscurity as readily as importance) and becomes a
question answered by a data test with a verdict a reader can inspect. That is a narrower claim than
``closing the LBD loop,'' and we prefer the narrower claim because it is the one the evidence supports.

Two properties of the loop do the work. The first is the willingness to answer \emph{no}, and to distinguish a
\emph{refuted} hypothesis from an \emph{untested} one. A generator with no discriminating filter turns one flood into
another; a filter that only ever confirms is not a filter. The knockdown-QC gate is the sharpest expression
of this --- a failed perturbation is reported as \emph{untested}, never as a false ``no effect'' --- and it is why the
approach can be trusted to reject as well as to nominate. The second is that the entire loop, generation
through self-audit, ran inside an agentic workbench in which an independent reviewer model enforced
calibrated language on the platform's own output. The paper's own thesis --- that machine-generated science
must carry receipts and honestly scoped verdicts --- was applied to the machine that generated it.

There is a replicability dividend here that we did not anticipate at the outset. The workbench exposes no
programmatic interface, so an analysis run in it is ordinarily a one-off, hand-paced web session. By driving
it headlessly and pulling its own provenance, we turned that one-off into a re-runnable, audited pipeline ---
and the sensitivity analyses in \S{}4.1b make the point concrete: the same panel, run locally and then inside
the workbench, reproduced byte-for-byte, verified by the workbench's own reviewer. An agentic analysis that
would otherwise be irreproducible-by-construction became replicable in principle --- subject to the
platform's user interface remaining stable (\S{}5.3). This is a small step, but a real one, toward
replicable-in-principle \emph{agentic} science rather than merely reproducible code.

\subsection*{5.2 The next experiment the approach nominates}\label{the-next-experiment-the-approach-nominates}

The value of a prioritization loop is not only that it culls, but that its survivors come with a sharp,
falsifiable next question. For NAB2, that question is directional. The Perturb-seq substrate tells us that
knocking NAB2 \emph{down} reshapes the Th1/Th2 program; it does not tell us whether, in disease, one would want
more NAB2 or less. We pursued that question orthogonally --- outside the Perturb-seq substrate --- and report the
result as a nomination for the next experiment, not as a finding.

Two lines of evidence bear on it. First, NAB2 is not a cancer dependency in the DepMap screens, which is
consistent with (though not proof of) a regulatory rather than an essential role --- non-contradictory context
rather than support. Second, and more informative, mining public expression data for the \emph{direction} of the
NAB2--atopic-eczema association: in lesional atopic-dermatitis skin, NAB2 is \textbf{down} relative to non-lesional
(bulk log\textsubscript{2}FC -0.32, FDR 0.002; anti-correlated with Th2 activity, $\rho$ -0.34), and a single-cell analysis
confirmed this is a per-cell reduction in the disease-relevant compartments (keratinocyte -0.51,
paired \emph{p} 0.027; T/NK -0.57) rather than a shift in cell-type composition. Read together, these are more
consistent with NAB2 acting as a \textbf{Th2 brake that is lost or suppressed in chronic lesions} than as a Th2
driver --- a reading that reconciles with the perturbation data (the brake is engaged during acute
polarization and lost in chronic disease; knocking it down removes it). If that reading holds, the naive
therapeutic move --- topical knockdown --- is likely \emph{backwards}, and the direction to test is
\textbf{restoration or up-modulation} of NAB2.

We are candid about the ceiling on this. Expression tracks the direction of \emph{association}, not the
distinction between an effector and a brake; the voting-arm evidence was mixed (one significant arm, others
ambiguous or null); and only a perturbation --- showing that \emph{restoring} NAB2 dampens the Th2 program --- can
convert this brake hypothesis into a directional therapeutic claim. That perturbation is the experiment the
loop nominates. This is exactly the shape we argue for: the approach does not end the inquiry, it points a
bench at the one experiment most worth running.

\subsection*{5.3 Limitations}\label{limitations}

We inherit the LBD tradition of a candid limitations section, and there is a great deal to be candid about.

\textbf{The verdicts are nominations, not causal claims.} Every ``supported'' verdict means the chain held with a
receipt at each hop, not that the biology is established. This is sharpest at the disease hop, which is an
\emph{association} receipt: the module--disease labels derive from genetic (GWAS-based) evidence with no
colocalization or linkage-disequilibrium control, so NAB2's atopic-eczema label could in principle be
LD-inherited from the adjacent 12q13 atopy locus. We falsified the \emph{cis}-effect on STAT6's expression, which
addresses the perturbation signal; it does not settle the provenance of the GWAS label, which we foreground
as the flagship nomination's key open question rather than treat as discharged (\S{}4.4b), and which only
variant-level colocalization against the 12q13 atopic-dermatitis GWAS can resolve.

\textbf{The disease hop's stringency is substrate-inherited, not the referee's own discrimination.} Our
label-shuffle control (\S{}4.1b) is the most important honesty check in the paper, and it returned a result we
report straight: the disease hop supports fewer than 1\% of arbitrary gene$\times$disease pairs, but so does a
random relabeling of the enrichment data --- the near-total refutation rate is a property of the enrichment
study's FDR family, inherited by the referee, not a discrimination the referee itself supplies. The
referee's confident-\emph{no} edge lies in the knockdown-QC gate and the effect/program demotions, which
are its own computations; the disease hop's contribution is a stringent, label-dependent nomination filter,
and we claim no more for it.

\textbf{The funnel numbers are conditioned, not free-standing.} Within the pre-gated universe the program hop is
a tautology (\texttt{refuted\_program\ $\equiv$\ 0}), so it discriminates in an individual receipt, not as a funnel filter;
and ``43 supported'' is a joint product of the literature-novelty gate and the referee (the referee alone
supports 395 of 47,220 pairs). We restate these in Results rather than burying them here.

\textbf{The substrate is bounded, and singular.} The prioritization is retrospective against one Perturb-seq
resource, one program (Th1/Th2), one cell type, and one set of aggregated tables; we ran no new wet-lab
experiment. The approach is only as good as the held substrate for a given claim, and it stays silent ---
returning \emph{untested} --- where the substrate cannot speak, which is a feature, but also a boundary on scope.

\textbf{The determinism boundary, and where model judgment enters.} Every data receipt is deterministic code; the
language model interprets receipts and never computes one. But the objective's weights are a human design
choice (we showed the headline is robust to them, \S{}4.1b), and the model's interpretive and provenance-writing
roles are judgment, not computation --- so the calibrated-language discipline, and the independent critic that
enforces it, are load-bearing rather than cosmetic.

\textbf{The self-audit is within-family, and the automation depends on an unstable surface.} The workbench's
reviewer is an independent \emph{role, model, and checkpoint}, but within a single model family; cross-family
independence is supplied externally (\S{}4.6), not natively. And the headless driver depends on the workbench's
web front-end, for which there is no stable public contract --- a replicability bridge over an API-less
platform, honestly, not a durable interface. If the front-end changes, the driver must be re-fitted.

\subsection*{5.4 Conclusion}\label{conclusion}

Literature-based discovery has never lacked for hypotheses; it has lacked a disciplined way to decide which
ones to pursue, and the candor to say \emph{no}. Wayfinder pairs a machine hypothesis-generator with a
deterministic referee that evaluates each proposal against a held experimental substrate and returns a
prioritization verdict with a receipt at every hop --- supporting, refuting, or declining to test --- and runs the whole loop,
including an audit of its own language, inside an agentic workbench whose work we can reproduce and inspect.
Applied to a genome-scale CD4\textsuperscript{+} T-cell resource, it culled tens of thousands of machine-generated hypotheses
to a small set of receipt-backed survivors, of which NAB2 $\rightarrow$ Th1/Th2 $\rightarrow$ atopic eczema is the worked example: a
literature-novel regulatory nomination whose sharpest artifactual confounder we falsified against the study
authors' own data, and which points to a concrete, falsifiable next experiment. The contribution is the
method and what it found --- receipt-backed prioritization of machine-generated hypotheses, with explicit
abstention and falsification diagnostics --- and the confident, receipt-backed \emph{no} that a triage step
must be willing to give.


\bibliography{../../../references/references}

\end{document}
